% Part: first-order-logic
% Chapter: sequent-calculus
% Section: soundness.tex

\documentclass[../../../include/open-logic-section]{subfiles}

\begin{document}

\iftag{FOL}
      {\olfileid[pt]{fol}{seq}{sou}}
      {\olfileid[pt]{pl}{seq}{sou}}

\olsection{Correção}

\begin{explain}
Um sistema de !!{derivation}, tal como o cálculo de sequentes, é
\emph{correto} se não pode derivar coisas que de fato não valem.  A
correção é, assim, uma espécie de propriedade de segurança garantida
para sistemas de !!{derivation}.  Dependendo de qual propriedade da
teoria da prova está em questão, gostaríamos de saber, por exemplo, que
\begin{enumerate}
\item todo~$!A$ !!{derivable} é \iftag{FOL}{válido}{uma tautologia};
\item se !!a{sentence} é !!{derivable} de algumas outras, ela também é
  consequência delas;
\item se um conjunto de !!{sentence}s é inconsistente, ele é insatisfatível.
\end{enumerate}
Essas são propriedades importantes de um sistema de !!{derivation}.  Se
alguma delas não vale, o sistema de !!{derivation} é deficiente---ele
!!{derive} demais.  Consequentemente, estabelecer a correção de um
sistema de !!{derivation} é da máxima importância.

Como todas essas propriedades da teoria da prova são definidas via
!!{derivability} no cálculo de sequentes de certos sequentes, provar
(1)--(3) acima requer provar algo sobre as propriedades semânticas dos
sequentes !!{derivable}s.  Primeiro definiremos o que significa um
sequente ser \emph{válido}, e então mostraremos que todo sequente
!!{derivable} é válido.  (1)--(3) seguem-se então como corolários
desse resultado.
\end{explain}

\begin{defn}
\iftag{FOL}{!!^a{structure}~$\Struct
  M$}{!!^a{valuation}~$\pAssign{v}$} \emph{satisfaz} um sequente
$\Gamma \Sequent \Delta$ se, e somente se, ou
$\iftag{FOL}{\Sat/{M}}{\pSat/{v}}{!A}$ para algum $!A \in \Gamma$ ou
$\iftag{FOL}{\Sat{M}}{\pSat{v}}{!A}$ para algum $!A \in \Delta$.

Um sequente é \emph{válido} se, e somente se, toda \iftag{FOL}{!!{structure}~$\Struct
  M$}{!!{valuation}~$\pAssign{v}$} o satisfaz.
\end{defn}

\begin{thm}[Correção]
\ollabel{thm:sequent-soundness} Se $\Log{LK}$ !!{derive} $\Theta
\Sequent \Xi$, então $\Theta \Sequent \Xi$ é válido.
\end{thm}

\begin{proof}
Seja $\pi$ !!a{derivation} de $\Theta \Sequent \Xi$. Procedemos por
indução sobre o número de inferências~$n$ em~$\pi$.

Se o número de inferências é~$0$, então $\pi$ consiste apenas em um
sequente inicial. Todo sequente inicial $!A \Sequent !A$ é obviamente
válido, pois para toda \iftag{FOL}{$\Struct M$}{$\pAssign{v}$}, ou
$\iftag{FOL}{\Sat/{M}}{\pSat/{v}}{!A}$ ou
$\iftag{FOL}{\Sat{M}}{\pSat{v}}{!A}$.

Se o número de inferências é maior que~0, distinguimos casos de acordo
com o tipo da inferência mais inferior. Pela hipótese de indução,
podemos supor que as premissas dessa inferência são válidas, já que o
número de inferências na !!{derivation} de qualquer premissa é menor
que~$n$.

Primeiro, consideramos as inferências possíveis com apenas uma premissa.
\begin{enumerate}
\item A última inferência é um enfraquecimento.  Então $\Theta \Sequent
  \Xi$ é $!A, \Gamma \Sequent \Delta$ (se a última inferência é
  \LeftR{\Weakening}) ou $\Gamma \Sequent \Delta, !A$ (se é
  \RightR{\Weakening}), e a !!{derivation} termina em uma de
  \begin{prooftree}
    \AxiomC{}
    \Deduce$\Gamma \fCenter \Delta$
    \RightLabel{\LeftR{\Weakening}}
    \UnaryInf$!A, \Gamma \fCenter \Delta$
    \DisplayProof\qquad\bottomAlignProof
    \AxiomC{}
    \Deduce$\Gamma \fCenter \Delta$
    \RightLabel{\RightR{\Weakening}}
    \UnaryInf$\Gamma \fCenter \Delta, !A$
  \end{prooftree}
  Pela hipótese de indução, $\Gamma \Sequent \Delta$ é válido, isto é,
  para toda
  \iftag{FOL}{!!{structure}~$\Struct{M}$}{!!{valuation}~$\pAssign{v}$},
  ou há algum $!C \in \Gamma$ tal que
  $\iftag{FOL}{\Sat/{M}}{\pSat/{v}}{!C}$ ou há algum $!C \in
  \Delta$ tal que $\iftag{FOL}{\Sat{M}}{\pSat{v}}{!C}$.
  
  Se $\iftag{FOL}{\Sat/{M}}{\pSat/{v}}{!C}$ para algum $!C \in \Gamma$,
  então $!C \in \Theta$ também, já que $\Theta = !A, \Gamma$, e assim
  $\iftag{FOL}{\Sat/{M}}{\pSat/{v}}{!C}$ para algum $!C \in \Theta$.
  Analogamente, se $\iftag{FOL}{\Sat{M}}{\pSat{v}}{!C}$ para algum $!C \in
  \Delta$, como $!C \in \Xi$, $\iftag{FOL}{\Sat{M}}{\pSat{v}}{!C}$ para
  algum $!C \in \Xi$. Consequentemente, $\Theta \Sequent \Xi$ é válido.
\item A última inferência é \LeftR{\lnot}: Então a premissa da última
  inferência é $\Gamma \Sequent \Delta, !A$ e a conclusão é
  $\lnot !A, \Gamma \Sequent \Delta$, isto é, a !!{derivation} termina em
  \begin{prooftree}
    \AxiomC{}
    \Deduce$\Gamma \fCenter \Delta, !A$
    \RightLabel{\LeftR{\lnot}}
    \UnaryInf$\lnot !A, \Gamma \fCenter \Delta$
  \end{prooftree}
  e $\Theta = \lnot !A, \Gamma$ enquanto $\Xi = \Delta$.

  A hipótese de indução nos diz que $\Gamma \Sequent \Delta, !A$
  é válido, isto é, para toda \iftag{FOL}{$\Struct{M}$}{$\pAssign{v}$},
  ou (a) para algum $!C \in \Gamma$,
  $\iftag{FOL}{\Sat/{M}}{\pSat/{v}}{!C}$, ou (b) para algum $!C \in
  \Delta$, $\iftag{FOL}{\Sat{M}}{\pSat{v}}{!C}$, ou (c)
  $\iftag{FOL}{\Sat{M}}{\pSat{v}}{!A}$. Queremos mostrar que $\Theta
  \Sequent \Xi$ também é válido.  Seja \iftag{FOL}{$\Struct{M}$
    !!a{structure}}{$\pAssign{v}$ !!a{valuation}}. Se (a) vale,
  então há $!C \in \Gamma$ tal que
  $\iftag{FOL}{\Sat/{M}}{\pSat/{v}}{!C}$, mas $!C \in \Theta$
  também. Se (b) vale, há $!C \in \Delta$ tal que
  $\iftag{FOL}{\Sat{M}}{\pSat{v}}{!C}$, mas $!C \in \Xi$
  também. Por fim, se $\iftag{FOL}{\Sat{M}}{\pSat{v}}{!A}$, então
  $\iftag{FOL}{\Sat/{M}}{\pSat/{v}}{\lnot !A}$. Como $\lnot !A \in
  \Theta$, há $!C \in \Theta$ tal que
  $\iftag{FOL}{\Sat/{M}}{\pSat/{v}}{!C}$. Consequentemente, $\Theta
  \Sequent \Xi$ é válido.
\item A última inferência é \RightR{\lnot}: Exercício.
\item A última inferência é \LeftR{\land}: Há duas variantes: $!A
  \land !B$ pode ser inferido à esquerda a partir de $!A$ ou de $!B$ no
  lado esquerdo da premissa.  No primeiro caso, $\pi$ termina em
  \begin{prooftree}
    \AxiomC{}
    \Deduce$!A, \Gamma \fCenter \Delta$
    \RightLabel{\LeftR{\land}}
    \UnaryInf$!A \land !B, \Gamma \fCenter \Delta$
  \end{prooftree}
  e $\Theta = !A \land !B, \Gamma$ enquanto $\Xi = \Delta$.  Considere
  \iftag{FOL}{!!a{structure}~$\Struct
    M$}{!!a{valuation}~$\pAssign{v}$}.  Como, pela hipótese de indução,
  $!A, \Gamma \Sequent \Delta$ é válido, (a)
  $\iftag{FOL}{\Sat/{M}}{\pSat/{v}}{!A}$, (b)
  $\iftag{FOL}{\Sat/{M}}{\pSat/{v}}{!C}$ para algum $!C \in \Gamma$, ou
  (c)~$\iftag{FOL}{\Sat{M}}{\pSat{v}}{!C}$ para algum $!C \in \Delta$.
  No caso (a), $\iftag{FOL}{\Sat/{M}}{\pSat/{v}}{!A \land !B}$, então
  há $!C \in \Theta$ (a saber, $!A \land !B$) tal que
  $\iftag{FOL}{\Sat/{M}}{\pSat/{v}}{!C}$.  No caso (b), há $!C
  \in \Gamma$ tal que $\iftag{FOL}{\Sat/{M}}{\pSat/{v}}{!C}$, e
  $!C \in \Theta$ também. No caso (c), há $!C \in \Delta$ tal
  que $\iftag{FOL}{\Sat{M}}{\pSat{v}}{!C}$, e $!C \in \Xi$ também,
  já que $\Xi = \Delta$. Assim, em cada caso, \iftag{FOL}{$\Struct
    M$}{$\pAssign{v}$} satisfaz $!A \land !B, \Gamma \Sequent
  \Delta$.  Como \iftag{FOL}{$\Struct{M}$}{$\pAssign{v}$} era
  arbitrário, $\Gamma \Sequent \Delta$ é válido.  O caso em que $!A
  \land !B$ é inferido a partir de $!B$ é tratado da mesma forma,
  mudando $!A$ para $!B$.
\item A última inferência é \RightR{\lor}: Há duas variantes: $!A
  \lor !B$ pode ser inferido à direita a partir de $!A$ ou de $!B$ no
  lado direito da premissa.  No primeiro caso, $\pi$ termina em
  \begin{prooftree}
    \AxiomC{}
    \Deduce$\Gamma \fCenter \Delta, !A$
    \RightLabel{\RightR{\lor}}
    \UnaryInf$\Gamma \fCenter \Delta, !A \lor !B$
  \end{prooftree}
  Agora $\Theta = \Gamma$ e $\Xi = \Delta, !A \lor !B$.  Considere
  \iftag{FOL}{!!a{structure}~$\Struct{M}$}{!!a{valuation}~$\pAssign{v}$}.
  Como $\Gamma \Sequent \Delta, !A$ é válido, (a)
  $\iftag{FOL}{\Sat{M}}{\pSat{v}}{!A}$, (b)
  $\iftag{FOL}{\Sat/{M}}{\pSat/{v}}{!C}$ para algum $!C \in \Gamma$, ou
  (c)~$\iftag{FOL}{\Sat{M}}{\pSat{v}}{!C}$ para algum $!C \in \Delta$.
  No caso (a), $\iftag{FOL}{\Sat{M}}{\pSat{v}}{!A \lor !B}$.  No caso
  (b), há $!C \in \Gamma$ tal que
  $\iftag{FOL}{\Sat/{M}}{\pSat/{v}}{!C}$.  No caso (c), há $!C
  \in \Delta$ tal que $\iftag{FOL}{\Sat{M}}{\pSat{v}}{!C}$. Assim, em
  cada caso, \iftag{FOL}{$\Struct{M}$}{$\pAssign{v}$} satisfaz
  $\Gamma \Sequent \Delta, !A \lor !B$, isto é, $\Theta \Sequent \Xi$.
  Como \iftag{FOL}{$\Struct{M}$}{$\pAssign{v}$} era arbitrário,
  $\Theta \Sequent \Xi$ é válido.  O caso em que $!A \lor !B$ é
  inferido a partir de $!B$ é tratado da mesma forma, mudando $!A$ para $!B$.
\item A última inferência é \RightR{\lif}: Então $\pi$ termina em
  \begin{prooftree}
    \AxiomC{}
    \Deduce$!A, \Gamma \fCenter \Delta, !B$
    \RightLabel{\RightR{\lif}}
    \UnaryInf$\Gamma \fCenter \Delta, !A \lif !B$
  \end{prooftree}
  Novamente, a hipótese de indução diz que a premissa é válida;
  queremos mostrar que a conclusão também é válida. Seja
  \iftag{FOL}{$\Struct{M}$}{$\pAssign{v}$} arbitrário. Como $!A,
  \Gamma \Sequent \Delta, !B$ é válido, ao menos um dos seguintes
  casos ocorre: (a) $\iftag{FOL}{\Sat/{M}}{\pSat/{v}}{!A}$, (b)
  $\iftag{FOL}{\Sat{M}}{\pSat{v}}{!B}$, (c)
  $\iftag{FOL}{\Sat/{M}}{\pSat/{v}}{!C}$ para algum~$~!C \in \Gamma$, ou
  (d) $\iftag{FOL}{\Sat{M}}{\pSat{v}}{!C}$ para algum $!C \in \Delta$.
  Nos casos (a) e (b), $\iftag{FOL}{\Sat{M}}{\pSat{v}}{!A \lif !B}$
  e assim há um $!C \in \Delta, !A \lif !B$ tal que
  $\iftag{FOL}{\Sat{M}}{\pSat{v}}{!C}$.  No caso (c), para algum $!C \in
  \Gamma$, $\iftag{FOL}{\Sat/{M}}{\pSat/{v}}{!C}$. No caso (d), para
  algum $!C \in \Delta$, $\iftag{FOL}{\Sat{M}}{\pSat{v}}{!C}$.  Em cada
  caso, \iftag{FOL}{$\Struct{M}$}{$\pAssign{v}$} satisfaz $\Gamma
  \Sequent \Delta, !A \lif !B$.  Como
  \iftag{FOL}{$\Struct{M}$}{$\pAssign{v}$} era arbitrário, $\Gamma
  \Sequent \Delta, !A \lif !B$ é válido.  \iftag{FOL}{%
\item A última inferência é \LeftR{\lforall}: Então há
  !!a{formula}~$!A(x)$ e um termo fechado~$t$ tais que $\pi$ termina em
  \begin{prooftree}
    \AxiomC{}
    \Deduce$!A(t), \Gamma \fCenter \Delta$
    \RightLabel{\LeftR{\lforall}}
    \UnaryInf$\lforall[x][!A(x)], \Gamma \fCenter \Delta$
  \end{prooftree}
  Queremos mostrar que a conclusão $\lforall[x][!A(x)], \Gamma
  \Sequent \Delta$ é válida.  Considere !!a{structure}~$\Struct M$.
  Como a premissa $!A(t), \Gamma \Sequent \Delta$ é válida, (a)
  $\Sat/{M}{!A(t)}$, (b) $\Sat/{M}{!C}$ para algum $!C \in \Gamma$, ou
  (c)~$\Sat{M}{!C}$ para algum $!C \in \Delta$.  No caso (a), por
  \olref[syn][sem]{prop:quant-terms}, se
  $\Sat{M}{\lforall[x][!A(x)]}$, então $\Sat{M}{!A(t)}$. Como
  $\Sat/{M}{!A(t)}$, $\Sat/{M}{\lforall[x][!A(x)]}$ . Nos casos (b) e
  (c), $\Struct{M}$ também satisfaz $\lforall[x][!A(x)], \Gamma
  \Sequent \Delta$.  Como $\Struct M$ era arbitrário,
  $\lforall[x][!A(x)], \Gamma \Sequent \Delta$ é válido.
\item A última inferência é \RightR{\lexists}: Exercício.
\item A última inferência é \RightR{\lforall}: Então há
  !!a{formula}~$!A(x)$ e !!a{constant}~$a$ tais que $\pi$ termina em
  \begin{prooftree}
    \AxiomC{}
    \Deduce$\Gamma \fCenter \Delta, !A(a)$
    \RightLabel{\RightR{\lforall}}
    \UnaryInf$\Gamma \fCenter \Delta, \lforall[x][!A(x)]$
  \end{prooftree}
  em que a condição de variável própria é satisfeita, isto é, $a$ não
  ocorre em $!A(x)$, $\Gamma$ ou $\Delta$.  Pela hipótese de indução,
  a premissa da última inferência é válida.  Temos de mostrar que
  a conclusão também é válida, isto é, que para qualquer
  !!{structure}~$\Struct M$, (a) $\Sat{M}{\lforall[x][!A(x)]}$, (b)
  $\Sat/{M}{!C}$ para algum $!C \in \Gamma$, ou (c)~$\Sat{M}{!C}$ para
  algum $!C \in \Delta$.

  Suponha que $\Struct{M}$ seja !!a{structure} arbitrária. Se (b) ou (c)
  vale, terminamos, então suponha que nenhuma vale: para todo $!C \in
  \Gamma$, $\Sat{M}{!C}$, e para todo $!C \in \Delta$,
  $\Sat/{M}{!C}$. Temos de mostrar que (a) vale, isto é,
  $\Sat{M}{\lforall[x][!A(x)]}$.  Por
  \olref[syn][ass]{prop:sat-quant}, basta mostrar que
  $\Sat{M}{!A(x)}[s]$ para toda atribuição de variáveis~$s$.  Então seja $s$
  uma atribuição de variáveis arbitrária. Considere a
  estrutura~$\Struct{M'}$ que é exatamente como~$\Struct{M}$ exceto que
  $\Assign{a}{M'} = s(x)$. Por
  \olref[syn][ext]{cor:extensionality-sent}, para qualquer $!C \in \Gamma$,
  $\Sat{M'}{!C}$ já que $a$ não ocorre em~$\Gamma$, e para qualquer
  $!C \in \Delta$, $\Sat/{M'}{!C}$. Mas a premissa é válida, então
  $\Sat{M'}{!A(a)}$. Por \olref[syn][ass]{prop:sentence-sat-true},
  $\Sat{M'}{!A(a)}[s]$, já que $!A(a)$ é uma sentença. Agora $\varAssign{s}{s}{x}$
  com $s(x) = \Value{a}{M'}[s]$, já que definimos $\Struct{M'}$
  exatamente desse modo. Assim \olref[syn][ext]{prop:ext-formulas} se aplica, e
  obtemos $\Sat{M'}{!A(x)}[s]$. Como $a$ não ocorre em~$!A(x)$, por
  \olref[syn][ext]{prop:extensionality}, $\Sat{M}{!A(x)}[s]$. Como
  $s$ era arbitrário, completamos a prova de que
  $\Sat{M}{!A(x)}[s]$ para toda atribuição de variáveis.
\item A última inferência é \LeftR{\lexists}: Exercício.
}{}
\end{enumerate}
Agora consideremos as inferências possíveis com duas premissas.
\begin{enumerate}
\item A última inferência é um corte: então $\pi$ termina em
  \begin{prooftree}
    \AxiomC{}
    \Deduce$\Gamma \fCenter \Delta, !A$
    \AxiomC{}
    \Deduce$!A, \Pi \fCenter \Lambda$
    \RightLabel{\Cut}
    \BinaryInf$\Gamma, \Pi \fCenter \Delta, \Lambda$
  \end{prooftree}
  Seja \iftag{FOL}{$\Struct{M}$ !!a{structure}}{$\pAssign{v}$
    !!a{valuation}}.  Pela hipótese de indução, as premissas são válidas,
  então \iftag{FOL}{$\Struct{M}$}{$\pAssign{v}$} satisfaz ambas as premissas.
  Distinguimos dois casos: (a) $\iftag{FOL}{\Sat/{M}}{\pSat/{v}}{!A}$
  e (b) $\iftag{FOL}{\Sat{M}}{\pSat{v}}{!A}$.  No caso (a), para que
  \iftag{FOL}{$\Struct{M}$}{$\pAssign{v}$} satisfaça a premissa
  esquerda, ela deve satisfazer $\Gamma \Sequent \Delta$.  Mas então
  também satisfaz a conclusão.  No caso (b), para que
  \iftag{FOL}{$\Struct{M}$}{$\pAssign{v}$} satisfaça a premissa
  direita, ela deve satisfazer $\Pi \setminus \Lambda$.  Novamente,
  \iftag{FOL}{$\Struct{M}$}{$\pAssign{v}$} satisfaz a conclusão.
\item A última inferência é \RightR{\land}.  Então $\pi$ termina em
  \begin{prooftree}
    \AxiomC{}
    \Deduce$\Gamma \fCenter \Delta, !A$
    \AxiomC{}
    \Deduce$\Gamma \fCenter \Delta, !B$
    \RightLabel{\RightR{\land}}
    \BinaryInf$\Gamma \fCenter \Delta, !A \land !B$
  \end{prooftree}
  Considere \iftag{FOL}{!!a{structure}~$\Struct
    M$}{!!a{valuation}~$\pAssign{v}$}. Se
  \iftag{FOL}{$\Struct{M}$}{$\pAssign{v}$} satisfaz $\Gamma \Sequent
  \Delta$, terminamos. Então suponha que não. Como $\Gamma \fCenter
  \Delta, !A$ é válido pela hipótese de indução,
  $\iftag{FOL}{\Sat{M}}{\pSat{v}}{!A}$. Analogamente, como $\Gamma
  \Sequent \Delta, !B$ é válido,
  $\iftag{FOL}{\Sat{M}}{\pSat{v}}{!B}$. Mas então
  $\iftag{FOL}{\Sat{M}}{\pSat{v}}{!A \land !B}$.
\item A última inferência é \LeftR{\lor}: Exercício.
\item A última inferência é \LeftR{\lif}.  Então $\pi$ termina em
  \begin{prooftree}
    \AxiomC{}
    \Deduce$\Gamma \fCenter \Delta, !A$
    \AxiomC{}
    \Deduce$!B, \Pi \fCenter \Lambda$
    \RightLabel{\LeftR{\lif}}
    \BinaryInf$!A \lif !B, \Gamma, \Pi \fCenter \Delta, \Lambda$
  \end{prooftree}
  Novamente, considere
  \iftag{FOL}{!!a{structure}~$\Struct{M}$}{!!a{valuation}~$\pAssign{v}$}
  e suponha que \iftag{FOL}{$\Struct{M}$}{$\pAssign{v}$} não satisfaça
  $\Gamma, \Pi \Sequent \Delta, \Lambda$. Temos de mostrar que
  $\iftag{FOL}{\Sat/{M}}{\pSat/{v}}{!A \lif !B}$. Se
  \iftag{FOL}{$\Struct{M}$}{$\pAssign{v}$} não satisfaz $\Gamma,
  \Pi \Sequent \Delta, \Lambda$, não satisfaz nem $\Gamma \Sequent
  \Delta$ nem $\Pi \Sequent \Lambda$. Como $\Gamma \Sequent \Delta,
  !A$ é válido, temos $\iftag{FOL}{\Sat{M}}{\pSat{v}}{!A}$. Como
  $!B, \Pi \Sequent \Lambda$ é válido, temos
  $\iftag{FOL}{\Sat/{M}}{\pSat/{v}}{!B}$. Mas então
  $\iftag{FOL}{\Sat/{M}}{\pSat/{v}}{!A \lif !B}$, que é o que
  queríamos mostrar.
\end{enumerate}
\end{proof}

\tagprob{FOL}
\begin{prob}
Complete a prova de \olref[fol][seq][sou]{thm:sequent-soundness}.
\end{prob}
\tagendprob

\tagprob{notFOL}
\begin{prob}
Complete a prova de \olref[pl][seq][sou]{thm:sequent-soundness}.
\end{prob}
\tagendprob

\begin{cor}
\ollabel{cor:weak-soundness}
Se $\Proves !A$, então $!A$ é \iftag{FOL}{válido}{uma tautologia}.
\end{cor}

\begin{cor}
\ollabel{cor:entailment-soundness}
Se $\Gamma \Proves !A$, então $\Gamma \Entails !A$.
\end{cor}

\begin{proof}
Se $\Gamma \Proves !A$, então, para algum subconjunto finito $\Gamma_0
\subseteq \Gamma$, há !!a{derivation} de $\Gamma_0 \Sequent !A$.  Por
\olref{thm:sequent-soundness}, toda
\iftag{FOL}{!!{structure}~$\Struct{M}$}{!!{valuation}~$\pAssign{v}$}
ou torna algum $!B \in \Gamma_0$ falso ou torna $!A$ verdadeiro.  Logo,
se $\iftag{FOL}{\Sat{M}}{\pSat{v}}{\Gamma}$, então também
$\iftag{FOL}{\Sat{M}}{\pSat{v}}{!A}$.
\end{proof}

\begin{cor}
\ollabel{cor:consistency-soundness}
Se $\Gamma$ é satisfatível, então é consistente.
\end{cor}

\begin{proof}
Provamos a contrapositiva.  Suponha que $\Gamma$ não seja consistente.
Então há um $\Gamma_0 \subseteq \Gamma$ finito e !!a{derivation}
de $\Gamma_0 \Sequent \quad$.  Por \olref{thm:sequent-soundness},
$\Gamma_0 \Sequent \quad$ é válido.  Em outras palavras, para toda
\iftag{FOL}{!!{structure}~$\Struct{M}$}{!!{valuation}~$\pAssign{v}$},
há $!C \in \Gamma_0$ tal que
$\iftag{FOL}{\Sat/{M}}{\pSat/{v}}{!C}$, e como $\Gamma_0
\subseteq \Gamma$, esse $!C$ também está em~$\Gamma$.  Assim, nenhuma
\iftag{FOL}{$\Struct{M}$}{$\pAssign{v}$} satisfaz~$\Gamma$, e
$\Gamma$ não é satisfatível.
\end{proof}

\end{document}
